% !TEX program = xelatex
% 编译方式：在 Overleaf 中选择 XeLaTeX，上传本文件和 latex_assets 文件夹。
% 本版说明书已依据 AttackPilot_backend 仓库结构重新核对。

\documentclass[UTF8,a4paper,12pt]{ctexart}

\usepackage[a4paper,left=22mm,right=22mm,top=23mm,bottom=22mm]{geometry}
\usepackage{graphicx}
\usepackage{booktabs}
\usepackage{tabularx}
\usepackage{array}
\usepackage{longtable}
\usepackage{enumitem}
\usepackage{float}
\usepackage{xcolor}
\usepackage{fancyhdr}
\usepackage{lastpage}
\usepackage{hyperref}
\usepackage{caption}

\definecolor{AttackBlue}{HTML}{1E4E79}
\definecolor{AttackGray}{HTML}{526173}

\hypersetup{
  colorlinks=true,
  linkcolor=AttackBlue,
  urlcolor=AttackBlue,
  citecolor=AttackBlue,
  pdftitle={AttackPilot作品设计说明书},
  pdfauthor={中山大学 AttackPilot 项目团队}
}

\setlength{\parindent}{2em}
\setlength{\parskip}{0.35em}
\setlength{\emergencystretch}{3em}
\renewcommand{\arraystretch}{1.35}
\captionsetup{font=small,labelfont=bf,justification=centering}

\pagestyle{fancy}
\fancyhf{}
\fancyhead[L]{\small AttackPilot 作品设计说明书}
\fancyhead[R]{\small 2026 年第三届全国大学生数智链应用大赛}
\fancyfoot[C]{\thepage\ / \pageref{LastPage}}
\renewcommand{\headrulewidth}{0.4pt}
\renewcommand{\footrulewidth}{0pt}

\newcommand{\asset}[2][0.88\textwidth]{%
  \begin{figure}[H]
    \centering
    \includegraphics[width=#1]{latex_assets/#2}
  \end{figure}
}

\newcommand{\todo}[1]{\textcolor{red}{\textbf{【提交前补填】#1}}}

\begin{document}

\begin{titlepage}
  \thispagestyle{empty}
  \vspace*{20mm}
  \begin{center}
    {\zihao{1}\bfseries\color{AttackBlue} AttackPilot 作品设计说明书\par}
    \vspace{8mm}
    {\zihao{3}\color{AttackGray}2026 年第三届全国大学生数智链应用大赛\par}
    \vspace{3mm}
    {\zihao{4}\color{AttackGray}区块链大类\quad 技术应用—应用实践与应用设计\par}
    \vspace{18mm}
    \renewcommand{\arraystretch}{1.7}
    \begin{tabularx}{0.84\textwidth}{>{\bfseries}p{0.23\textwidth}X}
      \toprule
      作品名称 & AttackPilot \\
      所选分类 & 技术应用—应用实践与应用设计 \\
      学校名称 & 中山大学 \\
      团队队长 & \todo{填写队长姓名} \\
      团队成员 & \todo{填写全部成员姓名} \\
      指导教师 & \todo{填写指导教师姓名} \\
      提交日期 & 2026 年 \todo{填写月份} 月 \todo{填写日期} 日 \\
      \bottomrule
    \end{tabularx}
    \vfill
    {\small 本说明书依据 AttackPilot 项目报告、作品材料和\texttt{AttackPilot\_backend}仓库代码核验整理。\par}
  \end{center}
\end{titlepage}

\tableofcontents
\clearpage

\section{作品概述}

\subsection{创作背景}

区块链账本公开、可追溯，但智能合约攻击通常不是一笔交易即可解释：攻击者可能先准备资金，再通过闪电贷、价格操纵、权限缺失、重入或错误会计等机制触发漏洞，最后完成资产转移。复盘人员需要同时查看交易收据、事件日志、内部调用、资金变化、合约源码和模拟结果，人工串联这些证据既耗时又容易遗漏。\par

AttackPilot 以链上攻击复盘为应用场景，把交易事实、Trace 执行证据、智能体分析、补丁编译与重放验证组织成一条可观察的工作流。作品的意义不是用模型替代安全人员，而是让模型分析始终有交易、源码和工具结果作为依据，并让最终结论可以追溯、复核和继续修正。

\subsection{作品简介}

AttackPilot 是一个面向 Ethereum 及 EVM 兼容链的链上攻击检测与复盘平台。用户可以从已整理的 DApp 案例进入，也可以提交一笔或多笔交易 Hash。后端先进行交易格式校验、案例匹配和快速风险筛选，再由任务管理器调度宏观交易分析、跨交易角色识别、故障定位、动态 Trace 调试、补丁生成、编译与模拟验证以及最终报告生成。\par

项目当前 GitHub 仓库为后端核心实现，提供 FastAPI REST API、WebSocket 实时通道、PostgreSQL/Redis 持久化、MCP 工具服务、Agent 运行链路、检查点恢复、进度 watchdog、结构化执行事件和本地 RAG 知识库。原作品材料中的前端工作台负责交易输入、任务控制台、日志、攻击路径和报告展示，完整参赛作品由前端展示层与本后端分析服务共同组成。

\subsection{目标受众与使用场景}

\begin{itemize}[leftmargin=2.5em]
  \item \textbf{区块链安全分析人员：}用于可疑交易初筛、攻击路径重建、证据核验和复盘报告整理。
  \item \textbf{智能合约开发者：}用于查看根因函数、触发条件、修复建议，并检查补丁能否阻断原始攻击。
  \item \textbf{教学与研究人员：}使用真实案例学习交易角色、漏洞机制、Trace 结构和修复方法。
  \item \textbf{审计与演示场景：}通过任务状态、事件链、检查点和案例审查结果展示一个可观察的智能合约分析过程。
\end{itemize}

\section{设计理念}

\subsection{设计主题与创意来源}

作品主题为“证据驱动的链上攻击复盘”。创意来自安全专家反复执行“提出假设—查找证据—验证结论—修正判断”的实际工作方式，也来自 TracePilot 跨交易故障定位研究。AttackPilot 将研究原型改造成可以接收任务、推送进度、保存中间产物、从检查点恢复并输出审查结果的应用服务。

\subsection{核心设计思路}

系统以任务为中心组织分析，不把一次模型回答当成最终结论。案件接入后，宏观分析阶段重建交易角色、资金变化和候选调试目标；任务规划阶段把宏观结论交接给下游 Agent；Transaction Debugger 使用 MCP 工具按需展开 Trace 节点、源码和状态；Code Patcher 生成补丁并编译；Transaction Judge 根据编译和模拟结果复核补丁；GlobalMemory Administrator 汇总跨交易上下文并生成报告。若证据不足、补丁失败或根因判断需要修正，系统记录真实的交接和反馈，而不是把流程强行渲染成一次成功的固定流水线。

\subsection{作品主要功能}

\begin{table}[H]
  \centering
  \caption{AttackPilot 主要功能与代码依据}
  \begin{tabularx}{0.96\textwidth}{>{\bfseries}p{0.23\textwidth}p{0.25\textwidth}X}
    \toprule
    功能模块 & 对外能力 & 后端实现依据 \\
    \midrule
    交易接入与检测 & 交易 Hash 校验、案例匹配、快速风险筛选、深度分析任务创建 & \texttt{/api/tx-detect}、\texttt{/api/tx-review}、\texttt{/api/tx-review/deep-analysis} \\
    多智能体复盘 & 宏观分析、角色识别、故障定位、Trace 调试、补丁与审查 & \texttt{main.py}、\texttt{process/}、\texttt{agents/}、\texttt{mcp\_tools/} \\
    实时任务工作台 & 任务状态、日志、进度、心跳和 WebSocket 推送 & \texttt{TaskManager}、\texttt{/ws/\{task\_id\}} \\
    可恢复执行 & 输入摘要、阶段检查点、SHA-256 完整性校验、恢复和取消 & \texttt{task\_checkpoint\_service.py}、\texttt{task\_progress\_service.py} \\
    结构化可观测性 & 工作流、阶段、Agent 交接、LLM 轮次和工具调用的父子事件链 & \texttt{execution-event-v1.schema.json}、\texttt{WorkflowEventRecorder} \\
    知识与审查 & 漏洞类型、精选案例、生成报告知识、混合检索、案例审查和质量门 & \texttt{data/knowledge/}、\texttt{data/rag/}、\texttt{case\_review\_service.py} \\
    \bottomrule
  \end{tabularx}
\end{table}

\section{开发过程}

\subsection{创作时间线}

项目报告描述了研究、平台开发、数据整理、测试和材料准备等阶段，但未提供完整的起止日期。正式提交前应根据 Git 记录、会议记录或实际开发日志补填下表；以下阶段划分与当前仓库的代码边界一致。

\begin{longtable}{>{\bfseries}p{0.22\textwidth}p{0.22\textwidth}p{0.47\textwidth}}
  \caption{AttackPilot 创作时间线}\label{tab:timeline}\\
  \toprule
  阶段 & 时间 & 主要工作 \\
  \midrule
  \endfirsthead
  \toprule
  阶段 & 时间 & 主要工作 \\
  \midrule
  \endhead
  需求与研究梳理 & \todo{起止日期} & 明确跨交易定位、动态 Trace、补丁验证和案例教学需求。 \\
  数据与分析原型 & \todo{起止日期} & 整理 DApp 攻击案例、交易数据、合约源码和 Trace 处理流程。 \\
  Agent 与 MCP 开发 & \todo{起止日期} & 实现宏观分析、任务规划、Trace 调试、补丁生成、审查和工具调用。 \\
  后端工程化 & \todo{起止日期} & 接入 FastAPI、任务调度、WebSocket、PostgreSQL、Redis 和 Docker Compose。 \\
  可靠性与审查增强 & \todo{起止日期} & 增加进度 watchdog、检查点恢复、结构化执行事件、RAG 和案例审查。 \\
  前端与竞赛材料 & \todo{起止日期} & 完成工作台、截图、视频、设计说明、部署说明和提交材料整理。 \\
  \bottomrule
\end{longtable}

\subsection{所使用的技术和系统架构}

AttackPilot_backend 是模块化单体后端。\texttt{app/main.py} 暴露 API 与 WebSocket；\texttt{TaskManager} 负责任务生命周期、线程池、取消、持久化和恢复；根目录\texttt{main.py} 的\texttt{WorkFlow} 负责编排实际分析；\texttt{process/} 将交易处理和 Agent 链路组合起来；\texttt{agents/} 实现各分析角色；\texttt{mcp\_tools/} 提供 Trace、源码、状态和补丁相关工具。\par

数据层使用 SQLAlchemy 对接 PostgreSQL，存储\texttt{task\_runs}、\texttt{task\_logs}、\texttt{task\_checkpoints}、\texttt{task\_progress\_events} 和\texttt{task\_execution\_events}；Redis 保存运行期高频上下文。\texttt{data/knowledge/} 保存漏洞类型、精选案例和生成报告知识，\texttt{data/rag/knowledge\_index.json} 保存本地检索索引，\texttt{dataset/raw/} 保存整理后的案例输入。Docker Compose 编排后端、PostgreSQL 和 Redis，默认通过 8000 端口提供服务。

\begin{table}[H]
  \centering
  \caption{当前仓库的系统分层}
  \begin{tabularx}{0.96\textwidth}{>{\bfseries}p{0.22\textwidth}p{0.27\textwidth}X}
    \toprule
    层次 & 主要组件 & 责任 \\
    \midrule
    接口层 & FastAPI、Uvicorn、Pydantic、WebSocket & 接收交易和 DApp 任务、返回报告、提供进度与事件查询、推送实时日志。 \\
    编排层 & TaskManager、WorkFlow、ThreadPool、Watchdog & 调度任务，维护状态，执行取消、检查点、恢复和进度策略。 \\
    智能体层 & TxDetail、TxRole、Filter、TxFault、Task、Trace、Fix、Judge、GlobalMemory & 完成交易详情、角色识别、攻击筛选、根因分析、Trace 调试、补丁与最终报告。 \\
    工具与取证层 & MCP Server、JSON-RPC、Etherscan 类接口、Tenderly、web3.py & 获取链上证据、读取源码、展开 Trace、编译和模拟补丁。 \\
    数据层 & PostgreSQL、Redis、JSON 案例、RAG 索引 & 保存任务、日志、检查点、事件、案例知识和运行时缓存。 \\
    \bottomrule
  \end{tabularx}
\end{table}

\asset{attackpilot-2.png}
\captionof{figure}{AttackPilot 整体技术架构示意图}

\subsection{接口与任务流程}

\begin{table}[H]
  \centering
  \caption{后端主要接口}
  \begin{tabularx}{0.96\textwidth}{p{0.12\textwidth}p{0.37\textwidth}X}
    \toprule
    方法 & 路径 & 说明 \\
    \midrule
    POST & \texttt{/api/tx-detect} & 对交易进行格式检查、链上观察、信号提取和案例匹配。 \\
    POST & \texttt{/api/tx-review/deep-analysis} & 根据交易 Hash 创建深度复盘任务。 \\
    POST/GET & \texttt{/api/tasks} & 创建任务、获取任务列表。 \\
    GET & \texttt{/api/tasks/\{task\_id\}/progress} & 查看阶段进度、健康状态、活动时间和 watchdog 信息。 \\
    GET/POST & \texttt{/api/tasks/\{task\_id\}/recovery}、\texttt{/resume} & 查看恢复计划并从有效检查点继续任务。 \\
    GET & \texttt{/api/tasks/\{task\_id\}/execution-events} & 分页读取结构化执行事件，可按需读取完整载荷。 \\
    GET & \texttt{/api/tasks/\{task\_id\}/case-review} & 返回带证据引用和信任等级的案例审查结果。 \\
    GET/POST & \texttt{/api/knowledge/*}、\texttt{/api/rag/*} & 浏览漏洞和案例知识、查询或重建 RAG 索引。 \\
    POST & \texttt{/api/assistant/chat} & 在任务、交易和知识库上下文中提供辅助问答。 \\
    WS & \texttt{/ws/\{task\_id\}} & 推送状态、日志、进度、检查点、心跳和最终结果。 \\
    \bottomrule
  \end{tabularx}
\end{table}

\section{区块链技术应用}

\subsection{使用的区块链平台和版本}

代码配置明确支持 Ethereum 主网和 BNBChain，并保留扩展到其他 EVM 网络的链 ID 映射。系统通过 JSON-RPC 获取交易、收据和 Trace，通过 Etherscan API v2 类接口获取合约源码和 ABI，通过 Tenderly 进行交易模拟、部署和补丁重放。Python 侧使用\texttt{web3.py}、\texttt{eth\_abi} 和\texttt{eth\_utils} 解析链上数据；Node 侧使用多版本\texttt{solc}、\texttt{solc-typed-ast} 和项目内的\texttt{solcjs.js}/\texttt{ast.js} 支持 Solidity 编译和源码映射。\par

项目不发行新代币、不修改公链共识，也不把私钥或真实 API 密钥作为作品功能。区块链在本作品中承担公开事实记录和智能合约执行证据来源，分析任务、日志和报告由应用层管理。

\begin{table}[H]
  \centering
  \caption{区块链相关技术及用途}
  \begin{tabularx}{0.96\textwidth}{p{0.23\textwidth}p{0.25\textwidth}X}
    \toprule
    平台或工具 & 代码依据 & 主要用途 \\
    \midrule
    Ethereum / BNBChain & \texttt{settings.py} 中的网络与链 ID 映射 & 获取真实交易、事件、Trace 和账户/合约信息。 \\
    JSON-RPC & \texttt{JSONRPCS}、\texttt{downloaders/} & 获取交易、收据、区块、调用和状态相关数据。 \\
    Etherscan API v2 类服务 & \texttt{SCAN\_APIKEYS}、源码下载器 & 读取验证源码、ABI、编译器配置和合约元数据。 \\
    Tenderly Simulation API & \texttt{daos/tenderly}、\texttt{FixAgent} & 部署或模拟补丁，并重放原始攻击路径。 \\
    Solidity 编译与 AST & \texttt{package.json}、\texttt{misc/ast.js}、\texttt{misc/solcjs.js} & 编译补丁、获得源码位置和执行节点映射。 \\
    \bottomrule
  \end{tabularx}
\end{table}

\subsection{区块链主要存储的信息类型}

平台读取和组织的区块链证据包括交易 Hash、区块号、交易收据、成功状态、调用者与被调用者、函数选择器、内部调用层级、事件日志、代币和原生币余额变化、合约源码、ABI、源码位置、状态读写以及模拟回放结果。分析流程将这些事实转换为攻击候选、交易角色、调试目标、根因假设、补丁结果和报告证据引用。\par

需要区分“链上事实”和“应用分析结果”：链上数据来自公链节点、扫描器或模拟器；任务状态、日志、检查点和报告存储在 PostgreSQL/Redis 或项目缓存中；模型生成的结论只有在关联到可核验证据后才进入案例审查结果。

\section{作品特色与创新}

\subsection{作品的创新点}

\begin{enumerate}[leftmargin=2.8em]
  \item \textbf{跨交易的攻击复盘：}以 DApp 案例为输入，联合分析准备、触发、套利和清理等交易，识别地址角色、资金变化和状态依赖。
  \item \textbf{动态 Trace 与 MCP 工具协作：}Agent 按当前假设选择展开节点、读取函数源码、查看状态和切换交易，避免一次性把完整 Trace 塞入模型上下文。
  \item \textbf{补丁生成到验证闭环：}FixAgent 读取缺陷合约和编译配置，生成补丁并尝试编译、部署或模拟；JudgeAgent 将结果反馈给后续判断。
  \item \textbf{可恢复的长任务执行：}检查点记录阶段、输入摘要、产物路径和 SHA-256，检测到有效产物后可从宏观分析、微观分析或报告恢复。
  \item \textbf{结构化执行证据链：}以 JSON Schema v1 记录\texttt{workflow.run}、\texttt{workflow.stage}、\texttt{agent.handoff}、\texttt{llm.turn} 和\texttt{tool.call}，用父事件和关联 ID 表达真实调用关系，并在写入前递归脱敏。
  \item \textbf{案例审查与知识回流：}将漏洞类型、精选案例、生成报告和检索结果统一为知识源；完成任务后可沉淀生成报告知识，案例审查服务再输出证据、完整性和信任等级。
\end{enumerate}

\subsection{与同类作品的区别与优势}

传统区块链工具通常只提供浏览器查询、Trace 展开、静态检测或单次模拟，用户仍需自行串联交易、源码、资金流和修复结果。AttackPilot 的差异在于以“任务”作为统一上下文，把交易接入、宏观筛选、Agent 协作、实时日志、执行事件、检查点、补丁验证和报告审查放在同一个可观测流程中。\par

与只保存文本日志的分析脚本相比，当前后端新增了结构化执行事件表和 JSON 契约，可回答“哪个 Agent 在哪个阶段调用了哪个工具、调用由什么事件触发、输出是否被修改、是否支撑了某个结论”等问题；与无恢复能力的长任务相比，当前后端以输入摘要和产物哈希避免错误复用；与只给出模型结论的系统相比，案例审查服务会区分模型报告、工具验证和证据完整性。

项目报告曾给出 149 个真实 DApp 安全事件的离线评测结果：AttackPilot 核心算法 Recall@Top-1、Recall@Top-3、Recall@Top-5 均为 71.14\%，Precision 为 77.78\%，平均端到端分析时间约 4,893.81 秒。该结果属于报告中的特定实验设置，不等同于当前仓库每次运行的实时保证；正式提交时应同时附上数据集、实验条件、脚本或可核验结果。

\begin{table}[H]
  \centering
  \caption{项目报告中的故障定位评测结果}
  \begin{tabular}{lcccc}
    \toprule
    方法 & Recall@1 & Recall@3 & Recall@5 & Precision \\
    \midrule
    DAppFL & 12.75\% & 19.46\% & 22.82\% & 4.56\% \\
    FaultSeeker & 31.08\% & 49.32\% & 50.67\% & 23.04\% \\
    AttackPilot 核心算法 & 71.14\% & 71.14\% & 71.14\% & 77.78\% \\
    \bottomrule
  \end{tabular}
\end{table}

\subsection{作品难点及解决方案}

\begin{enumerate}[leftmargin=2.8em]
  \item \textbf{跨交易关系复杂：}利用交易角色、资金变化、函数入口、调用层级和事件信息建立上下文，再选择深度调试交易。
  \item \textbf{Trace 体量大：}通过 MCP 按需展开节点，TraceAgent 根据上下文预算裁剪内容，并保留关键节点、源码和状态读写。
  \item \textbf{Agent 过程难以解释：}新增追加式执行事件，使用\texttt{parent\_event\_id}、\texttt{correlation\_id}、预览、摘要和证据引用构造调用链。
  \item \textbf{任务耗时长且可能中断：}采用阶段检查点、产物完整性校验、进度百分比和 watchdog；日志活动不被误当作真实进度。
  \item \textbf{补丁可能无法编译或无法阻断攻击：}将编译、模拟和回放作为质量门，区分模型报告的“已验证”与具备工具产物的“工具验证”。
  \item \textbf{外部服务不稳定：}对 RPC、扫描器、Tenderly 和 LLM 调用采用配置化、超时、并发桶和失败反馈，并保留离线案例、RAG 索引和 smoke test。
\end{enumerate}

\section{创作总结}

\subsection{开发过程中的收获与经验}

项目实践表明，模型能力只是智能安全应用的一部分。真正影响可用性的还有数据结构、任务生命周期、网络超时、日志语义、产物完整性和用户是否能看懂证据。把研究流程拆成输入、宏观分析、任务规划、Trace 调试、补丁验证和报告阶段，有助于逐步检查质量，也便于前端向用户解释当前任务处于哪里。

\subsection{区块链使用的体会与反思}

区块链数据公开且不可随意改写，但交易事实的理解仍需要上下文。单独看一笔交易或一条事件日志，可能无法解释攻击者如何准备状态、何时改变价格、谁最终获得资产。系统应把链上数据作为事实来源，把模型输出作为待验证假设，把源码、Trace、模拟和回放作为证据。对于 RPC、扫描器和模拟服务，还要同时考虑接口配额、网络失败、密钥权限和结果差异。

\subsection{作品的不足之处及改进方向}

当前版本仍依赖外部 RPC、扫描器、Tenderly 和 OpenAI-compatible LLM 服务；完整新案例分析需要网络和相应配置，演示仍应准备经过脱敏的缓存案例。当前仓库中的前端不在本次克隆的后端目录内，正式参赛包必须将前端、后端、安装说明、演示视频和部署地址对应好。\par

后续可改进方向包括：增加更多 EVM 网络适配；替换仓库中可能存在的硬编码服务凭据并完善密钥管理；补充 CI、数据库迁移和容器健康检查；把更多 MCP 工具结果类型化；增加人工确认和误报分析；提供可复现实验脚本、统一报告导出和在线部署监控。

\section{参考资料}

\subsection{设计开发过程中参考的文献和资料}

\begin{enumerate}[leftmargin=2.8em,label={[\arabic*]}]
  \item AttackPilot 项目报告书，团队内部项目材料。
  \item AttackPilot Backend，\url{https://github.com/Tlipped/AttackPilot_backend}。
  \item TracePilot: Self-verifiable Framework for Decentralized Applications Fault Localization across Transactions，ISSTA 2026。
  \item Sun, K.; Xu, Z.; Li, K.; Zhang, L.; Sun, Y.; Tan, L.; Liu, Y. FAULTSEEKER: LLM-Empowered Framework for Blockchain Transaction Fault Localization. ASE 2025。
  \item Wu, Z.; Wu, J.; Zhang, H.; Li, Z.; Chen, J.; Zheng, Z.; Xia, Q.; Fan, G.; Zhen, Y. DAppFL: Just-in-Time Fault Localization for Decentralized Applications in Web3. ISSTA 2024, 137--148。
  \item SunWeb3Sec. DeFiHackLabs. \url{https://github.com/SunWeb3Sec/DeFiHackLabs}。
  \item BlockSec. Phalcon Explorer. \url{https://blocksec.com/phalcon/explorer}。
  \item Tenderly Developer Explorer and Simulation API. \url{https://tenderly.co/developer-explorer}。
  \item Solidity Documentation. \url{https://docs.soliditylang.org/}。
  \item Ethereum Yellow Paper，G. Wood，2014。
\end{enumerate}

\subsection{第三方资源使用情况说明及授权依据}

后端使用 FastAPI、Uvicorn、Pydantic、WebSocket、SQLAlchemy、PostgreSQL、Redis、Docker Compose、web3.py、eth-abi、eth-utils、FastMCP、MCP、OpenAI Python SDK、Transformers、NetworkX、Matplotlib 等开源组件，分别用于接口服务、数据校验、实时通信、持久化、缓存、容器编排、链上数据处理、工具调用和分析辅助。Node 依赖包括多版本 Solidity 编译器和\texttt{solc-typed-ast}。具体版本以仓库中的\texttt{requirements.txt}、\texttt{package.json} 和 Docker 配置为准。\par

项目参考了 DeFiHackLabs、Phalcon、Tenderly、Ethernaut 等公开案例、工具或教学资源。提交前应逐项核对许可证、版权声明和引用方式，保留第三方原始声明，不将外部案例或库误写为团队原创。AI 工具仅用于资料阅读、代码辅助、文档整理和运行问题排查；问题定义、数据核对、系统设计、测试判断和最终提交由团队成员确认并负责。

\section*{附录}
\addcontentsline{toc}{section}{附录}

\subsection*{附录 A：作品成品展示截图与效果图}
\addcontentsline{toc}{subsection}{附录 A：作品成品展示截图与效果图}

以下图片来自原作品材料，用于展示完整作品的前端工作台与设计方案；当前 GitHub 克隆仓库主要承担其后端服务。

\asset{attackpilot-1.png}
\captionof{figure}{AttackPilot 产品首页}
\asset{attackpilot-3.png}
\captionof{figure}{用户功能模块}
\asset{attackpilot-4.png}
\captionof{figure}{多智能体系统流程}
\asset{attackpilot-5.png}
\captionof{figure}{项目报告中的时间效率对比}
\asset{attackpilot-6.png}
\captionof{figure}{任务输入工作区}
\asset{attackpilot-7.png}
\captionof{figure}{漏洞教学区}
\asset{attackpilot-8.png}
\captionof{figure}{分析复盘区}

\subsection*{附录 B：区块链配置信息或区块链浏览器截图}
\addcontentsline{toc}{subsection}{附录 B：区块链配置信息或区块链浏览器截图}

正式提交时应补充至少一组脱敏材料：
\begin{enumerate}[leftmargin=2.8em]
  \item 链上交易浏览器页面，能看到交易 Hash、区块、状态、调用或事件日志；
  \item 后端运行配置截图，仅保留网络名称、接口类型和服务状态，不显示 API Key、Authorization、私钥、密码或完整访问令牌；
  \item 如展示 Tenderly 或其他模拟服务，应遮挡账户标识、项目密钥和可复用的访问链接。
\end{enumerate}
\todo{将脱敏截图保存为图片后插入本附录。}

\subsection*{附录 C：部署与运行说明}
\addcontentsline{toc}{subsection}{附录 C：部署与运行说明}

\textbf{后端仓库：} \url{https://github.com/Tlipped/AttackPilot_backend}。本说明只描述当前后端仓库，前端工作台需按其独立仓库或作品材料的说明启动。

\begin{enumerate}[leftmargin=2.8em]
  \item 复制\texttt{.env\_example}为\texttt{.env}，填写模型服务配置；如需实时链上分析，再配置 Tenderly、RPC 和扫描器服务。
  \item 在仓库根目录执行\texttt{docker compose up -d --build}，Docker Compose 会启动后端、PostgreSQL 和 Redis。
  \item 打开\url{http://localhost:8000/docs} 查看 Swagger 接口文档；用\url{http://localhost:8000/api/tasks} 检查任务 API。
  \item 使用\texttt{docker compose logs -f backend} 查看后端日志，使用\texttt{docker compose down} 停止服务。
  \item 在不访问真实链上网络的条件下，可运行\texttt{python scripts/smoke\_backend.py} 进行关键文件编译、Hash 归一化和本地 RAG 检查。
\end{enumerate}

当前仓库包含\texttt{docs/execution-event-contract.md}、\texttt{docs/workflow-handoff-contract.md} 和\texttt{docs/schemas/execution-event-v1.schema.json}，可作为结构化执行事件和 Agent 交接的技术补充材料。提交源代码前必须检查\texttt{settings.py}、\texttt{.env} 和历史提交，删除或替换任何真实服务凭据；作品包只保留脱敏配置和必要的示例值。

\subsection*{附录 D：仓库核验清单}
\addcontentsline{toc}{subsection}{附录 D：仓库核验清单}

\begin{table}[H]
  \centering
  \caption{本版说明书对应的仓库核验范围}
  \begin{tabularx}{0.96\textwidth}{>{\bfseries}p{0.30\textwidth}p{0.18\textwidth}X}
    \toprule
    核验对象 & 当前状态 & 用途 \\
    \midrule
    FastAPI 与 WebSocket 接口 & 已发现 & 任务、检测、审查、知识、RAG、事件和实时推送。 \\
    PostgreSQL/Redis 数据层 & 已发现 & 任务、原始日志、检查点、进度事件、结构化执行事件和缓存。 \\
    多智能体与 MCP & 已发现 & 宏观分析、Trace 调试、补丁生成、验证和报告。 \\
    检查点与 watchdog & 已发现 & 长任务进度、停滞判断、取消和安全恢复。 \\
    执行事件契约与测试 & 已发现 & 结构化调用链、脱敏、摘要、完整性和 API 回归测试。 \\
    前端展示、演示视频和公网部署 & 以原作品材料为准 & 需要在竞赛提交包中与本后端版本对应并补齐链接。 \\
    团队信息、区块链截图、真实部署地址 & 待补填 & 不应由代码仓库推断，必须由团队确认。 \\
    \bottomrule
  \end{tabularx}
\end{table}

\end{document}
